% GENERATED FILE. Edit canonical lessons, then run npm run generate:books.
% canonical-source-hash: fnv1a64:ce37947193bced6e
% canonical-lessons: ML-C205-lalitamaya, ML-C205-upayogapradamaya, ML-C205-mikacca, ML-C205-vicitramaya, ML-C205-adhunika

\chapter{Simple, Useful, Excellent, Strange, Modern}
\label{ch:ml-simple-useful-excellent-strange-modern}
\hlchaptermodality{205}

\begin{chapteropening}
\textbf{By the end of this chapter:} \emph{I can say simple, useful, excellent, strange and modern.}

Complete the last lesson of chapter 205: I can say simple, useful, excellent, strange and modern.
\end{chapteropening}

\section[\texorpdfstring{\ml{ലളിതമായ} (laḷitamāya)}{ലളിതമായ (laḷitamāya)}]{\ml{ലളിതമായ} (laḷitamāya) — simple}

\label{lesson:ML-C205-lalitamaya}



\begin{quote}
Before the new one: say the Malayalam for rough, then the Malayalam for round.
\end{quote}

\subsection*{You'll want to know: \ml{ലളിതമായ}}
\textbf{\ml{ലളിതമായ}} — \emph{laḷitamāya} — \textquotedblleft{}simple\textquotedblright{}.

\begin{cousinweb}
\textbf{\ml{ലളിതമായ}} (\emph{laḷitamāya}) joins Sanskrit \emph{lalita}, \textquotedblleft{}graceful, easy\textquotedblright{}, to native \emph{āya}, \textquotedblleft{}having become\textquotedblright{} (DEDR 333). DEDR 333 compares Sanskrit \emph{ām}.
\end{cousinweb}

\begin{grammarlens}[title={What you've built}]
One more word for this chapter.
\end{grammarlens}

\subsection*{Your turn}
\begin{itemize}
\raggedright
  \item \emph{Say it:} \emph{laḷitamāya}
  \item \emph{Say it:} \emph{laḷitamāya}, once more
  \item \emph{Say it:} say \emph{laḷitamāya}
  \item \emph{Recall:} say the Malayalam for rough, then the Malayalam for round, then say \emph{laḷitamāya} again
\end{itemize}

\begin{tcolorbox}[breakable,colback=teal!4,colframe=teal!35!black,title={Before you move on}]
What does \ml{ലളിതമായ} mean? (\textquotedblleft{}Simple\textquotedblright{}.) Say it once more.
\end{tcolorbox}
\section[\texorpdfstring{\ml{ഉപയോഗപ്രദമായ} (upayōgapradamāya)}{ഉപയോഗപ്രദമായ (upayōgapradamāya)}]{\ml{ഉപയോഗപ്രദമായ} (upayōgapradamāya) — useful}

\label{lesson:ML-C205-upayogapradamaya}



\begin{quote}
Before the new one: say the Malayalam for round, then the Malayalam for simple.
\end{quote}

\subsection*{You'll want to know: \ml{ഉപയോഗപ്രദമായ}}
\textbf{\ml{ഉപയോഗപ്രദമായ}} — \emph{upayōgapradamāya} — \textquotedblleft{}useful\textquotedblright{}.

\begin{cousinweb}
\textbf{\ml{ഉപയോഗപ്രദമായ}} (\emph{upayōgapradamāya}) is Sanskrit \emph{upayōga}, \textquotedblleft{}use\textquotedblright{}, and \emph{prada}, \textquotedblleft{}giving\textquotedblright{}, with native \emph{āya}, \textquotedblleft{}having become\textquotedblright{} (DEDR 333): \textquotedblleft{}use-giving\textquotedblright{}. DEDR 333 compares Sanskrit \emph{ām}.
\end{cousinweb}

\begin{grammarlens}[title={What you've built}]
One more word for this chapter.
\end{grammarlens}

\subsection*{Your turn}
\begin{itemize}
\raggedright
  \item \emph{Say it:} \emph{upayōgapradamāya}
  \item \emph{Say it:} \emph{upayōgapradamāya}, once more
  \item \emph{Say it:} say \emph{upayōgapradamāya}
  \item \emph{Recall:} say the Malayalam for round, then the Malayalam for simple, then say \emph{upayōgapradamāya} again
\end{itemize}

\begin{tcolorbox}[breakable,colback=teal!4,colframe=teal!35!black,title={Before you move on}]
What does \ml{ഉപയോഗപ്രദമായ} mean? (\textquotedblleft{}Useful\textquotedblright{}.) Say it once more.
\end{tcolorbox}
\section[\texorpdfstring{\ml{മികച്ച} (mikacca)}{മികച്ച (mikacca)}]{\ml{മികച്ച} (mikacca) — excellent, best}

\label{lesson:ML-C205-mikacca}



\begin{quote}
Before the new one: say the Malayalam for simple, then the Malayalam for useful.
\end{quote}

\subsection*{You'll want to know: \ml{മികച്ച}}
\textbf{\ml{മികച്ച}} — \emph{mikacca} — \textquotedblleft{}excellent, best\textquotedblright{}.

\begin{cousinweb}
\textbf{\ml{മികച്ച}} (\emph{mikacca}) is native Dravidian, a past form of the verb \textquotedblleft{}to increase, be foremost\textquotedblright{}; DEDR 4838 lists the root with Tamil \emph{miku}, Kannada \emph{mikku} and Telugu \emph{migulu}.
\end{cousinweb}

\begin{grammarlens}[title={What you've built}]
One more word for this chapter.
\end{grammarlens}

\subsection*{Your turn}
\begin{itemize}
\raggedright
  \item \emph{Say it:} \emph{mikacca}
  \item \emph{Say it:} \emph{mikacca}, once more
  \item \emph{Say it:} say \emph{mikacca}
  \item \emph{Recall:} say the Malayalam for simple, then the Malayalam for useful, then say \emph{mikacca} again
\end{itemize}

\begin{tcolorbox}[breakable,colback=teal!4,colframe=teal!35!black,title={Before you move on}]
What does \ml{മികച്ച} mean? (\textquotedblleft{}Excellent, best\textquotedblright{}.) Say it once more.
\end{tcolorbox}
\section[\texorpdfstring{\ml{വിചിത്രമായ} (vicitramāya)}{വിചിത്രമായ (vicitramāya)}]{\ml{വിചിത്രമായ} (vicitramāya) — strange, odd}

\label{lesson:ML-C205-vicitramaya}



\begin{quote}
Before the new one: say the Malayalam for useful, then the Malayalam for excellent.
\end{quote}

\subsection*{You'll want to know: \ml{വിചിത്രമായ}}
\textbf{\ml{വിചിത്രമായ}} — \emph{vicitramāya} — \textquotedblleft{}strange, odd\textquotedblright{}.

\begin{cousinweb}
\textbf{\ml{വിചിത്രമായ}} (\emph{vicitramāya}) joins Sanskrit \emph{vicitra}, \textquotedblleft{}many-coloured, strange\textquotedblright{}, to native \emph{āya}, \textquotedblleft{}having become\textquotedblright{} (DEDR 333). DEDR 333 compares Sanskrit \emph{ām}.
\end{cousinweb}

\begin{grammarlens}[title={What you've built}]
One more word for this chapter.
\end{grammarlens}

\subsection*{Your turn}
\begin{itemize}
\raggedright
  \item \emph{Say it:} \emph{vicitramāya}
  \item \emph{Say it:} \emph{vicitramāya}, once more
  \item \emph{Say it:} say \emph{vicitramāya}
  \item \emph{Recall:} say the Malayalam for useful, then the Malayalam for excellent, then say \emph{vicitramāya} again
\end{itemize}

\begin{tcolorbox}[breakable,colback=teal!4,colframe=teal!35!black,title={Before you move on}]
What does \ml{വിചിത്രമായ} mean? (\textquotedblleft{}Strange, odd\textquotedblright{}.) Say it once more.
\end{tcolorbox}
\section[\texorpdfstring{\ml{ആധുനിക} (ādhunika)}{ആധുനിക (ādhunika)}]{\ml{ആധുനിക} (ādhunika) — modern}

\label{lesson:ML-C205-adhunika}



\begin{quote}
Before the new one: say the Malayalam for excellent, then the Malayalam for strange.
\end{quote}

\subsection*{You'll want to know: \ml{ആധുനിക}}
\textbf{\ml{ആധുനിക}} — \emph{ādhunika} — \textquotedblleft{}modern\textquotedblright{}.

\begin{cousinweb}
\textbf{\ml{ആധുനിക}} (\emph{ādhunika}) $\leftarrow$ Sanskrit \emph{ādhunika}, \textquotedblleft{}recent, new\textquotedblright{}, taken over whole as a tatsama. It is built on \emph{adhunā}, \textquotedblleft{}now\textquotedblright{}.
\end{cousinweb}

\begin{grammarlens}[title={What you've built}]
One more word for this chapter.
\end{grammarlens}

\subsection*{Your turn}
\begin{itemize}
\raggedright
  \item \emph{Say it:} \emph{ādhunika}
  \item \emph{Say it:} \emph{ādhunika}, once more
  \item \emph{Say it:} say \emph{ādhunika}
  \item \emph{Recall:} say the Malayalam for excellent, then the Malayalam for strange, then say \emph{ādhunika} again
\end{itemize}

\begin{tcolorbox}[breakable,colback=teal!4,colframe=teal!35!black,title={Before you move on}]
What does \ml{ആധുനിക} mean? (\textquotedblleft{}Modern\textquotedblright{}.) Say it once more.
\end{tcolorbox}
